\section{Compressing a graphing after enlarging a supplied relation}
\label{compression:section}

This section proves a cost estimate for replacing part of a graphing by a
larger relation whose edges are supplied without charge.  The measure on
the underlying space is finite and need not be a probability measure.

Let $(Z,\lambda)$ be a standard Borel space with a finite Borel measure.
Let $\mathcal R$ be a Borel equivalence relation covered by the graphs of
countably many measure-preserving Borel partial isomorphisms
$(\psi_j)_{j\geq 1}$.  A partial isomorphism means a Borel bijection between
Borel subsets with Borel inverse.  For a Borel subrelation $\mathcal U$ of
$\mathcal R$ and a Borel set $W\subseteq Z$, put
\[
 \kappa(\mathcal U;W)
 =\int_W\frac{1}{|[z]_{\mathcal U}|}\,d\lambda(z),
 \qquad
 \kappa(\mathcal U)=\kappa(\mathcal U;Z),
\]
where the reciprocal of an infinite class size is zero.  A set is
$\mathcal U$-complete if it meets every $\mathcal U$-class.
All relations in this section are pointwise Borel relations.  When an
application starts with identities holding almost everywhere, it may
first remove the $\mathcal R$-saturation of the exceptional null set.
That saturation is Borel and null because the maps $\psi_j$ are
measure-preserving.

\subsection{Borel selections and small complete sets}

We record the elementary measurability facts used in the construction.
Restricting $\psi_j$ to the set where
$(z,\psi_j(z))\in\mathcal U$ gives an enumeration of any Borel
subrelation $\mathcal U$.  Its saturations of Borel sets are therefore
countable unions of Borel images and are Borel.  Its class-size function
is Borel as well: having at least $k$ class members is a countable union
over $k$-tuples of enumeration values that are defined and pairwise
distinct.

Every Borel partial isomorphism $h$ with graph in $\mathcal R$ preserves
measure.  Indeed, partition its domain according to the first index $j$
for which $h(z)=\psi_j(z)$.  Each restriction preserves measure, and the
images of the pieces are disjoint because $h$ is injective.

Fix a Borel linear order on $Z$, obtained from a Borel injection into the
real line.  Nonempty finite subsets have least points.  For a Borel
relation with finite classes, selecting this least point is Borel: among
the countably enumerated candidates, the condition that none is smaller
is Borel.  The same observation applies to nonempty Borel subsets of
finite classes and to the successive points in their ordered lists.

\begin{lemma}[Finite-class counting]\label{compression:counting}
Let $W\subseteq Z$ be a Borel $\mathcal U$-invariant set all of whose
$\mathcal U$-classes are finite.  If $B\subseteq W$ is Borel and meets
each such class exactly once, then
\[
 \lambda(B)=\kappa(\mathcal U;W).
\]
\end{lemma}

\begin{proof}
On the part $W_k$ consisting of classes of size $k$, order each class
into $k$ Borel sheets.  The map from $B\cap W_k$ to any one sheet that
stays within each class is a Borel partial isomorphism inside
$\mathcal R$, so all $k$ sheets have measure $\lambda(B\cap W_k)$.
Thus $\lambda(W_k)=k\lambda(B\cap W_k)$.  Summing over $k$ proves the
identity.
\end{proof}

\begin{lemma}[Small complete sets]\label{compression:complete-set}
For every Borel subrelation $\mathcal T\subseteq\mathcal R$ and every
$\varepsilon>0$, there is a Borel $\mathcal T$-complete set
$M\subseteq Z$ such that
\[
 \lambda(M)\leq\kappa(\mathcal T)+\varepsilon.
\]
\end{lemma}

\begin{proof}
On the union $Z_{\mathrm{fin}}$ of finite $\mathcal T$-classes, choose
the least point of each class.  The resulting Borel set
$M_{\mathrm{fin}}$ has measure $\kappa(\mathcal T)$ by
Lemma~\ref{compression:counting}.

On $Z_{\mathrm{inf}}=Z\setminus Z_{\mathrm{fin}}$, choose increasing
finite Borel partitions $\mathcal P_n$ that separate all points, for
example those generated by the first $n$ members of a countable
separating Borel family.  Let $k_n(z)$ be the number of atoms of
$\mathcal P_n$ meeting $[z]_{\mathcal T}$.  This function is Borel, and
$k_n(z)\longrightarrow\infty$: any prescribed finite collection of
distinct points in the infinite class is separated by some partition.

Fix $0<p<1$ with $p\lambda(Z_{\mathrm{inf}})<\varepsilon/2$.
At a fixed level $n$, mark each atom independently with probability $p$,
and let $A$ be the union of the marked atoms.  The Borel set
\[
 M_{\mathrm{inf}}
 =A\cup\bigl(Z_{\mathrm{inf}}\setminus[A]_{\mathcal T}\bigr)
\]
meets every infinite $\mathcal T$-class for every marking.  The displayed
union is disjoint, and
\[
 \mathbb E\lambda(M_{\mathrm{inf}})
 =p\lambda(Z_{\mathrm{inf}})
   +\int_{Z_{\mathrm{inf}}}(1-p)^{k_n(z)}\,d\lambda(z).
\]
The integral tends to zero by dominated convergence.  For sufficiently
large $n$ the expectation is less than $\varepsilon$, so one of the
finitely many markings has measure less than $\varepsilon$.  Fix that
marking and set $M=M_{\mathrm{fin}}\cup M_{\mathrm{inf}}$.
\end{proof}

\subsection{Compression and its cost}

For a countable family $\mathcal E=(\varphi_i:D_i\to E_i)$ of
measure-preserving Borel partial isomorphisms, write
$C(\mathcal E)=\sum_i\lambda(D_i)$.  The join
$\mathcal S\vee\mathcal E$ denotes the equivalence relation generated
by $\mathcal S$ and all pairs $(x,\varphi_i(x))$.

\begin{lemma}[Compression]\label{compression:main}\label{compression:lemma}
On the finite-measure space $(Z,\lambda)$ with the ambient relation
$\mathcal R$ described above, suppose
$\mathcal S\subseteq\mathcal T\subseteq\mathcal Q\subseteq
\mathcal R$ are Borel equivalence relations and
$\mathcal Q=\mathcal S\vee\mathcal E$, where $\mathcal E$ is a
countable family of measure-preserving Borel partial isomorphisms with
finite total cost.  For every $\varepsilon>0$ there is a countable
family $\mathcal F$ of measure-preserving Borel partial isomorphisms
inside $\mathcal Q$ such that
\begin{equation}\label{compression:bound}
 \mathcal Q=\mathcal T\vee\mathcal F,
 \qquad
 C(\mathcal F)\leq C(\mathcal E)-\kappa(\mathcal S)
                          +\kappa(\mathcal T)+\varepsilon.
\end{equation}
\end{lemma}

\begin{proof}
Choose a Borel $\mathcal T$-complete set $M$ from
Lemma~\ref{compression:complete-set}, and put
\[
 Y=[M]_{\mathcal S}
     \cup\{z:|[z]_{\mathcal S}|=\infty\}.
\]
This set is Borel and $\mathcal S$-invariant.  It is
$\mathcal T$-complete and hence $\mathcal Q$-complete, since it contains
$M$.  Every $\mathcal S$-class outside $Y$ is finite.  In each finite
$\mathcal S$-class in $Y$, choose a point of its intersection with $M$.
Lemma~\ref{compression:counting} gives
\begin{equation}\label{compression:retained-classes}
 \kappa(\mathcal S;Y)\leq\lambda(M).
\end{equation}

\paragraph{Choosing and deleting one edge per class outside $Y$.}
Let $d(z)$ be the least number of edges from $\mathcal E$, in either
direction, needed to reach $Y$, allowing $\mathcal S$-steps for free.
This distance is finite because $Y$ is $\mathcal Q$-complete.  Its
sublevel sets are Borel: starting with $Y$, successively take images
under the maps of $\mathcal E$ and their inverses, and take
$\mathcal S$-saturations.  It is constant on $\mathcal S$-classes and
vanishes precisely on $Y$.

For every $\mathcal S$-class $K$ outside $Y$, choose an oriented edge
instance from $y_K\in K$ to $y'_K$ with $d(y'_K)=d(K)-1$.
Such an instance exists by a shortest path.  To make the choice Borel,
first choose the least eligible pair consisting of a map index and a
direction, then choose the least eligible point in the finite class
$K$.  Eligibility of a map index and direction is detected by the
$\mathcal S$-saturation of a Borel set.  Thus all these choices are
Borel.  The set $V=\{y_K\}$ is a Borel transversal of
$\mathcal S$ on $Z\setminus Y$, so
\begin{equation}\label{compression:selector-cost}
 \lambda(V)=\kappa(\mathcal S;Z\setminus Y).
\end{equation}

For each original map $\varphi_i:D_i\to E_i$, let $V_i^+$ be the
chosen origins whose edge is traversed forward under $\varphi_i$, and
let $V_i^-$ be those whose edge is traversed under $\varphi_i^{-1}$.
Delete the domain set
\[
 H_i=V_i^+\cup\varphi_i^{-1}(V_i^-)
\]
from $D_i$.  The union is disjoint: a labeled edge instance chosen in
both directions would make each endpoint class have smaller distance
than the other.  The sets $V_i^+,V_i^-$ partition $V$, and inverse
images under $\varphi_i$ preserve measure.  Consequently the residual
family $\mathcal E'=(\varphi_i|_{D_i\setminus H_i})$ satisfies the
exact identity
\begin{equation}\label{compression:deleted-cost}
 C(\mathcal E')
 =C(\mathcal E)-\sum_i\lambda(H_i)
 =C(\mathcal E)-\kappa(\mathcal S;Z\setminus Y).
\end{equation}
Different original indices count as different edge instances, even
when their geometric edges coincide.  Thus this accounting also covers
graphings with repetitions.

\paragraph{Moving the remaining edges into $Y$.}
Define $f:Z\to Y$ by $f(x)=x$ for $x\in Y$, and, for $x$ in an
$\mathcal S$-class $K$ outside $Y$, by
\[
 f(x)=f(y'_K).
\]
The recursion decreases the integer distance at every step.  Defining
$f$ successively on the distance levels proves it is Borel.  Each
$(x,f(x))$ lies in $\mathcal Q$.  Because $Y$ is
$\mathcal S$-invariant, $f$ fixes each class inside $Y$ pointwise and
is constant on each class outside $Y$.  It therefore sends
$\mathcal S$-steps to $\mathcal S$-steps.  Each deleted edge has equal
images at its two endpoints.

The map $f$ can be many-to-one.  To obtain partial isomorphisms after
moving the remaining edges, we must split their domains at both
endpoints.  Fix a residual edge map $\varphi_i$ with domain $D'_i$.
Partition $D'_i$ into Borel pieces $P_{i,r,s}$ according to the least
indices $r,s$ for which
\[
 f(x)=\psi_r(x),\qquad
 f(\varphi_i(x))=\psi_s(\varphi_i(x)).
\]
These indices exist because the graph of $f$ is contained in
$\mathcal Q\subseteq\mathcal R$.  On this piece the pushed edge is
the partial isomorphism
\[
 \widehat\varphi_{i,r,s}
 =\psi_s\circ\varphi_i\circ\psi_r^{-1}:
 \psi_r(P_{i,r,s})\longrightarrow
 \psi_s(\varphi_i(P_{i,r,s})).
\]
Its domain and range are Borel, it preserves measure, and its graph is
contained in $\mathcal Q|Y$.  Its cost equals
$\lambda(P_{i,r,s})$.  Let $\mathcal F$ consist of all these maps.
Although the images of distinct pieces may overlap, graphing cost is
the sum over the indexed family, and the original pieces partition
each $D'_i$.  Hence
\begin{equation}\label{compression:pushed-cost}
 C(\mathcal F)=\sum_{i,r,s}\lambda(P_{i,r,s})=C(\mathcal E').
\end{equation}
One may discard identity maps or redundant instances, in which case
the equality becomes an inequality in the needed direction.

\paragraph{Recovering the relation.}
For two $\mathcal Q$-related points in $Y$, take a finite path with
steps in $\mathcal S$ and $\mathcal E$ and apply $f$ to its vertices.
The endpoints remain fixed.  Each $\mathcal S$-step remains such a
step, each deleted edge becomes stationary, and each surviving edge
becomes an edge of $\mathcal F$.  Therefore
\[
 \mathcal Q|Y=(\mathcal S|Y)\vee\mathcal F.
\]
For arbitrary $x\mathrel{\mathcal Q}x'$, choose
$y\in[x]_{\mathcal T}\cap Y$ and
$y'\in[x']_{\mathcal T}\cap Y$ using completeness of $Y$.
Then $y\mathrel{\mathcal Q}y'$, so the preceding identity joins them
using $\mathcal S\subseteq\mathcal T$ and $\mathcal F$.
Together with the first and last $\mathcal T$-steps this proves
$\mathcal Q=\mathcal T\vee\mathcal F$.

Finally, by \eqref{compression:retained-classes},
\eqref{compression:deleted-cost}, and
\eqref{compression:pushed-cost},
\[
 \begin{aligned}
 C(\mathcal F)
 &\leq C(\mathcal E)-\kappa(\mathcal S;Z\setminus Y)\\
 &=C(\mathcal E)-\kappa(\mathcal S)+\kappa(\mathcal S;Y)\\
 &\leq C(\mathcal E)-\kappa(\mathcal S)+\lambda(M)\\
 &\leq C(\mathcal E)-\kappa(\mathcal S)
                       +\kappa(\mathcal T)+\varepsilon.
 \end{aligned}
\]
This is \eqref{compression:bound}.
\end{proof}

\begin{corollary}\label{compression:cost-one}
A countable p.m.p. Borel equivalence relation on a standard probability
space, all of whose classes are infinite, has cost at least one.
\end{corollary}
\begin{proof}
Let $\mathcal E$ be a finite-cost generating graphing of the relation
$\mathcal Q$. Apply Lemma~\ref{compression:main} with
$\mathcal S=\mathrm{id}$ and $\mathcal T=\mathcal Q$. Since
$\kappa(\mathcal S)=1$ and $\kappa(\mathcal Q)=0$, it gives
$0\le C(\mathcal F)\le C(\mathcal E)-1+\varepsilon$ for every
$\varepsilon>0$. Thus $C(\mathcal E)\ge1$. Infinite-cost graphings
satisfy this bound as well; taking the infimum proves the assertion.
\end{proof}

\begin{remark}\label{compression:normalization}
The estimate uses the original finite measure throughout; in
particular $\kappa(\mathrm{id}_Z)=\lambda(Z)$.  The output graphing
stays inside the same relation $\mathcal Q$, but no inclusion between
the original and replacement edge families is asserted.  These are
the two features needed when applying the lemma repeatedly on an
enlarged space.
\end{remark}
